\documentclass[10pt,letterpaper]{article}
\usepackage{cornell-notes}
\usepackage{mathtools}

\title{Chapter 6: Special Functions}
\author{}
\date{}

\setDocTitle{Chapter 6: Special Functions}
\setDocSubtitle{Numerical Methods}
\setDocVersion{v1.0}
\setDocStatus{Study Companion}
\setDocOwner{Numerical Methods Cornell Notes}
\setDocAuthor{}
\setDocDate{\today}
\setDocSummary{Cornell-format study and review notes for Chapter 6: Special Functions.}

\setCornellCollection{Numerical Methods Cornell Notes}
\setCornellUnitType{Chapter}
\setCornellUnitNumber{6}
\setCornellUnitTitle{Special Functions}
\setCornellCompanionLabel{Study and review companion}

\hypersetup{pdftitle={Chapter 6: Special Functions},pdfsubject={Cornell notes},pdfauthor={}}

\begin{document}
\maketitle

\section*{Chapter Overview}
This chapter organizes the mathematical and computational ideas behind special functions. The notes preserve the numbered section sequence while emphasizing method selection, explicit assumptions, numerical reliability, cost, and verification.

\section*{Learning Objectives}
Explain the governing problem class; compare Gamma Function, Beta Function, Factorials, Binomial Coefficients, Incomplete Gamma Function and Error Function, Exponential Integrals, and Incomplete Beta Function; design defensible stopping and verification checks; distinguish problem conditioning from algorithmic stability.

\section*{Prerequisites}
Calculus, differential equations, series, asymptotics, and complex arithmetic.

\section*{Operational Workflow}
Classify the argument region and requested scaling; use identities to reach a stable region; choose series, recurrence, continued fraction, asymptotic expansion, or path continuation; cross-check identities.

\subsection*{Formula and notation anchor}
\[
\Gamma(z)=\int_0^\infty t^{z-1}e^{-t}\,dt,\qquad B(a,b)=\frac{\Gamma(a)\Gamma(b)}{\Gamma(a+b)}
\]
Log-gamma evaluation is usually required for large arguments; the integral definition alone is not a general computational algorithm.

\begin{warnbox}{Numerical warning}
Overflow before cancellation, wrong recurrence direction, branch errors, loss of tail probability, discontinuous region switching, and inaccurate normalization.
\end{warnbox}

\section{6.0 Introduction}
\begin{cornell}
\noterow{What is the central idea?}{Special-function software partitions the domain and chooses among series, asymptotics, recurrences, continued fractions, and transformations.}
\noterow{How should it be used?}{For Introduction, begin from explicit inputs, outputs, preconditions, and representation choices.  Classify the argument region and requested scaling; use identities to reach a stable region; choose series, recurrence, continued fraction, asymptotic expansion, or path continuation; cross-check identities.}
\noterow{Cost and reliability}{Analyze arithmetic work, storage, data movement, and convergence separately.  Relevant failure pressures include overflow before cancellation, wrong recurrence direction, branch errors, loss of tail probability, discontinuous region switching, and inaccurate normalization.  Use scaled tolerances and report both success criteria and diagnostic evidence.}
\noterow{Implementation checklist}{\begin{itemize}
\item Validate dimensions, domains, ordering, and structural assumptions.
\item Guard small divisors, invalid states, overflow, underflow, and nonfinite values.
\item Make mutation, workspace, indexing, and termination behavior explicit.
\item Test a simple exact case, a difficult valid case, a boundary case, and an expected failure.
\end{itemize}}
\end{cornell}
\begin{summarybox}{Section summary}
Special-function software partitions the domain and chooses among series, asymptotics, recurrences, continued fractions, and transformations.  Select this technique only when its assumptions match the data, and accept a result only after an independent residual, error estimate, invariant, or refinement check.
\end{summarybox}

\section{6.1 Gamma Function, Beta Function, Factorials, Binomial Coefficients}
\begin{cornell}
\noterow{What is the central idea?}{Logarithmic forms and identities extend factorial-like quantities while avoiding immediate overflow and preserving sign information.}
\noterow{How should it be used?}{For Gamma Function, Beta Function, Factorials, Binomial Coefficients, begin from explicit inputs, outputs, preconditions, and representation choices.  Classify the argument region and requested scaling; use identities to reach a stable region; choose series, recurrence, continued fraction, asymptotic expansion, or path continuation; cross-check identities.}
\noterow{Quantitative anchor}{\[
\Gamma(z+1)=z\Gamma(z)
\]
State the dimensions, scaling, and validity assumptions before using this relation.  Check the resulting residual or invariant rather than trusting an iteration count alone.}
\noterow{Cost and reliability}{Analyze arithmetic work, storage, data movement, and convergence separately.  Relevant failure pressures include overflow before cancellation, wrong recurrence direction, branch errors, loss of tail probability, discontinuous region switching, and inaccurate normalization.  Use scaled tolerances and report both success criteria and diagnostic evidence.}
\noterow{Implementation checklist}{\begin{itemize}
\item Validate dimensions, domains, ordering, and structural assumptions.
\item Guard small divisors, invalid states, overflow, underflow, and nonfinite values.
\item Make mutation, workspace, indexing, and termination behavior explicit.
\item Test a simple exact case, a difficult valid case, a boundary case, and an expected failure.
\end{itemize}}
\end{cornell}
\begin{summarybox}{Section summary}
Logarithmic forms and identities extend factorial-like quantities while avoiding immediate overflow and preserving sign information.  Select this technique only when its assumptions match the data, and accept a result only after an independent residual, error estimate, invariant, or refinement check.
\end{summarybox}

\section{6.2 Incomplete Gamma Function and Error Function}
\begin{cornell}
\noterow{What is the central idea?}{Series and continued fractions cover complementary regions; selecting by argument prevents slow convergence and loss of accuracy.}
\noterow{How should it be used?}{For Incomplete Gamma Function and Error Function, begin from explicit inputs, outputs, preconditions, and representation choices.  Classify the argument region and requested scaling; use identities to reach a stable region; choose series, recurrence, continued fraction, asymptotic expansion, or path continuation; cross-check identities.}
\noterow{Quantitative anchor}{\[
\Gamma(z+1)=z\Gamma(z)
\]
State the dimensions, scaling, and validity assumptions before using this relation.  Check the resulting residual or invariant rather than trusting an iteration count alone.}
\noterow{Cost and reliability}{Analyze arithmetic work, storage, data movement, and convergence separately.  Relevant failure pressures include overflow before cancellation, wrong recurrence direction, branch errors, loss of tail probability, discontinuous region switching, and inaccurate normalization.  Use scaled tolerances and report both success criteria and diagnostic evidence.}
\noterow{Implementation checklist}{\begin{itemize}
\item Validate dimensions, domains, ordering, and structural assumptions.
\item Guard small divisors, invalid states, overflow, underflow, and nonfinite values.
\item Make mutation, workspace, indexing, and termination behavior explicit.
\item Test a simple exact case, a difficult valid case, a boundary case, and an expected failure.
\end{itemize}}
\end{cornell}
\begin{summarybox}{Section summary}
Series and continued fractions cover complementary regions; selecting by argument prevents slow convergence and loss of accuracy.  Select this technique only when its assumptions match the data, and accept a result only after an independent residual, error estimate, invariant, or refinement check.
\end{summarybox}

\section{6.3 Exponential Integrals}
\begin{cornell}
\noterow{What is the central idea?}{Exponential integrals need different representations near singular points, across branch-sensitive regions, and for large arguments.}
\noterow{How should it be used?}{For Exponential Integrals, begin from explicit inputs, outputs, preconditions, and representation choices.  Classify the argument region and requested scaling; use identities to reach a stable region; choose series, recurrence, continued fraction, asymptotic expansion, or path continuation; cross-check identities.}
\noterow{Quantitative lens}{Identify the governing equation, recurrence, probability law, or geometric predicate for Exponential Integrals.  Define every variable and scale; then derive a residual, error estimate, or invariant that can be checked independently.}
\noterow{Cost and reliability}{Analyze arithmetic work, storage, data movement, and convergence separately.  Relevant failure pressures include overflow before cancellation, wrong recurrence direction, branch errors, loss of tail probability, discontinuous region switching, and inaccurate normalization.  Use scaled tolerances and report both success criteria and diagnostic evidence.}
\noterow{Implementation checklist}{\begin{itemize}
\item Validate dimensions, domains, ordering, and structural assumptions.
\item Guard small divisors, invalid states, overflow, underflow, and nonfinite values.
\item Make mutation, workspace, indexing, and termination behavior explicit.
\item Test a simple exact case, a difficult valid case, a boundary case, and an expected failure.
\end{itemize}}
\end{cornell}
\begin{summarybox}{Section summary}
Exponential integrals need different representations near singular points, across branch-sensitive regions, and for large arguments.  Select this technique only when its assumptions match the data, and accept a result only after an independent residual, error estimate, invariant, or refinement check.
\end{summarybox}

\section{6.4 Incomplete Beta Function}
\begin{cornell}
\noterow{What is the central idea?}{Symmetry and continued-fraction evaluation move the calculation into a well-conditioned tail and support cumulative probability functions.}
\noterow{How should it be used?}{For Incomplete Beta Function, begin from explicit inputs, outputs, preconditions, and representation choices.  Classify the argument region and requested scaling; use identities to reach a stable region; choose series, recurrence, continued fraction, asymptotic expansion, or path continuation; cross-check identities.}
\noterow{Quantitative lens}{Identify the governing equation, recurrence, probability law, or geometric predicate for Incomplete Beta Function.  Define every variable and scale; then derive a residual, error estimate, or invariant that can be checked independently.}
\noterow{Cost and reliability}{Analyze arithmetic work, storage, data movement, and convergence separately.  Relevant failure pressures include overflow before cancellation, wrong recurrence direction, branch errors, loss of tail probability, discontinuous region switching, and inaccurate normalization.  Use scaled tolerances and report both success criteria and diagnostic evidence.}
\noterow{Implementation checklist}{\begin{itemize}
\item Validate dimensions, domains, ordering, and structural assumptions.
\item Guard small divisors, invalid states, overflow, underflow, and nonfinite values.
\item Make mutation, workspace, indexing, and termination behavior explicit.
\item Test a simple exact case, a difficult valid case, a boundary case, and an expected failure.
\end{itemize}}
\end{cornell}
\begin{summarybox}{Section summary}
Symmetry and continued-fraction evaluation move the calculation into a well-conditioned tail and support cumulative probability functions.  Select this technique only when its assumptions match the data, and accept a result only after an independent residual, error estimate, invariant, or refinement check.
\end{summarybox}

\section{6.5 Bessel Functions of Integer Order}
\begin{cornell}
\noterow{What is the central idea?}{Bessel recurrences, asymptotic forms, and normalization must be combined in the stable direction for the requested order and argument.}
\noterow{How should it be used?}{For Bessel Functions of Integer Order, begin from explicit inputs, outputs, preconditions, and representation choices.  Classify the argument region and requested scaling; use identities to reach a stable region; choose series, recurrence, continued fraction, asymptotic expansion, or path continuation; cross-check identities.}
\noterow{Quantitative lens}{Identify the governing equation, recurrence, probability law, or geometric predicate for Bessel Functions of Integer Order.  Define every variable and scale; then derive a residual, error estimate, or invariant that can be checked independently.}
\noterow{Cost and reliability}{Analyze arithmetic work, storage, data movement, and convergence separately.  Relevant failure pressures include overflow before cancellation, wrong recurrence direction, branch errors, loss of tail probability, discontinuous region switching, and inaccurate normalization.  Use scaled tolerances and report both success criteria and diagnostic evidence.}
\noterow{Implementation checklist}{\begin{itemize}
\item Validate dimensions, domains, ordering, and structural assumptions.
\item Guard small divisors, invalid states, overflow, underflow, and nonfinite values.
\item Make mutation, workspace, indexing, and termination behavior explicit.
\item Test a simple exact case, a difficult valid case, a boundary case, and an expected failure.
\end{itemize}}
\end{cornell}
\begin{summarybox}{Section summary}
Bessel recurrences, asymptotic forms, and normalization must be combined in the stable direction for the requested order and argument.  Select this technique only when its assumptions match the data, and accept a result only after an independent residual, error estimate, invariant, or refinement check.
\end{summarybox}

\section{6.6 Bessel Functions of Fractional Order, Airy Functions, Spherical Bessel Functions}
\begin{cornell}
\noterow{What is the central idea?}{Related function families can be reduced to shared base evaluations, but scale and phase conventions must remain consistent.}
\noterow{How should it be used?}{For Bessel Functions of Fractional Order, Airy Functions, Spherical Bessel Functions, begin from explicit inputs, outputs, preconditions, and representation choices.  Classify the argument region and requested scaling; use identities to reach a stable region; choose series, recurrence, continued fraction, asymptotic expansion, or path continuation; cross-check identities.}
\noterow{Quantitative lens}{Identify the governing equation, recurrence, probability law, or geometric predicate for Bessel Functions of Fractional Order, Airy Functions, Spherical Bessel Functions.  Define every variable and scale; then derive a residual, error estimate, or invariant that can be checked independently.}
\noterow{Cost and reliability}{Analyze arithmetic work, storage, data movement, and convergence separately.  Relevant failure pressures include overflow before cancellation, wrong recurrence direction, branch errors, loss of tail probability, discontinuous region switching, and inaccurate normalization.  Use scaled tolerances and report both success criteria and diagnostic evidence.}
\noterow{Implementation checklist}{\begin{itemize}
\item Validate dimensions, domains, ordering, and structural assumptions.
\item Guard small divisors, invalid states, overflow, underflow, and nonfinite values.
\item Make mutation, workspace, indexing, and termination behavior explicit.
\item Test a simple exact case, a difficult valid case, a boundary case, and an expected failure.
\end{itemize}}
\end{cornell}
\begin{summarybox}{Section summary}
Related function families can be reduced to shared base evaluations, but scale and phase conventions must remain consistent.  Select this technique only when its assumptions match the data, and accept a result only after an independent residual, error estimate, invariant, or refinement check.
\end{summarybox}

\section{6.7 Spherical Harmonics}
\begin{cornell}
\noterow{What is the central idea?}{Spherical harmonics combine angular exponentials with associated Legendre functions to form an orthogonal basis on the sphere.}
\noterow{How should it be used?}{For Spherical Harmonics, begin from explicit inputs, outputs, preconditions, and representation choices.  Classify the argument region and requested scaling; use identities to reach a stable region; choose series, recurrence, continued fraction, asymptotic expansion, or path continuation; cross-check identities.}
\noterow{Quantitative lens}{Identify the governing equation, recurrence, probability law, or geometric predicate for Spherical Harmonics.  Define every variable and scale; then derive a residual, error estimate, or invariant that can be checked independently.}
\noterow{Cost and reliability}{Analyze arithmetic work, storage, data movement, and convergence separately.  Relevant failure pressures include overflow before cancellation, wrong recurrence direction, branch errors, loss of tail probability, discontinuous region switching, and inaccurate normalization.  Use scaled tolerances and report both success criteria and diagnostic evidence.}
\noterow{Implementation checklist}{\begin{itemize}
\item Validate dimensions, domains, ordering, and structural assumptions.
\item Guard small divisors, invalid states, overflow, underflow, and nonfinite values.
\item Make mutation, workspace, indexing, and termination behavior explicit.
\item Test a simple exact case, a difficult valid case, a boundary case, and an expected failure.
\end{itemize}}
\end{cornell}
\begin{summarybox}{Section summary}
Spherical harmonics combine angular exponentials with associated Legendre functions to form an orthogonal basis on the sphere.  Select this technique only when its assumptions match the data, and accept a result only after an independent residual, error estimate, invariant, or refinement check.
\end{summarybox}

\section{6.8 Fresnel Integrals, Cosine and Sine Integrals}
\begin{cornell}
\noterow{What is the central idea?}{Oscillatory integrals demand region-dependent series, rational approximations, or asymptotic expansions to avoid inefficient direct quadrature.}
\noterow{How should it be used?}{For Fresnel Integrals, Cosine and Sine Integrals, begin from explicit inputs, outputs, preconditions, and representation choices.  Classify the argument region and requested scaling; use identities to reach a stable region; choose series, recurrence, continued fraction, asymptotic expansion, or path continuation; cross-check identities.}
\noterow{Quantitative lens}{Identify the governing equation, recurrence, probability law, or geometric predicate for Fresnel Integrals, Cosine and Sine Integrals.  Define every variable and scale; then derive a residual, error estimate, or invariant that can be checked independently.}
\noterow{Cost and reliability}{Analyze arithmetic work, storage, data movement, and convergence separately.  Relevant failure pressures include overflow before cancellation, wrong recurrence direction, branch errors, loss of tail probability, discontinuous region switching, and inaccurate normalization.  Use scaled tolerances and report both success criteria and diagnostic evidence.}
\noterow{Implementation checklist}{\begin{itemize}
\item Validate dimensions, domains, ordering, and structural assumptions.
\item Guard small divisors, invalid states, overflow, underflow, and nonfinite values.
\item Make mutation, workspace, indexing, and termination behavior explicit.
\item Test a simple exact case, a difficult valid case, a boundary case, and an expected failure.
\end{itemize}}
\end{cornell}
\begin{summarybox}{Section summary}
Oscillatory integrals demand region-dependent series, rational approximations, or asymptotic expansions to avoid inefficient direct quadrature.  Select this technique only when its assumptions match the data, and accept a result only after an independent residual, error estimate, invariant, or refinement check.
\end{summarybox}

\section{6.9 Dawson's Integral}
\begin{cornell}
\noterow{What is the central idea?}{Dawson's integral admits stable local series and asymptotic behavior, providing a scaled alternative to exponentially large intermediate expressions.}
\noterow{How should it be used?}{For Dawson's Integral, begin from explicit inputs, outputs, preconditions, and representation choices.  Classify the argument region and requested scaling; use identities to reach a stable region; choose series, recurrence, continued fraction, asymptotic expansion, or path continuation; cross-check identities.}
\noterow{Quantitative lens}{Identify the governing equation, recurrence, probability law, or geometric predicate for Dawson's Integral.  Define every variable and scale; then derive a residual, error estimate, or invariant that can be checked independently.}
\noterow{Cost and reliability}{Analyze arithmetic work, storage, data movement, and convergence separately.  Relevant failure pressures include overflow before cancellation, wrong recurrence direction, branch errors, loss of tail probability, discontinuous region switching, and inaccurate normalization.  Use scaled tolerances and report both success criteria and diagnostic evidence.}
\noterow{Implementation checklist}{\begin{itemize}
\item Validate dimensions, domains, ordering, and structural assumptions.
\item Guard small divisors, invalid states, overflow, underflow, and nonfinite values.
\item Make mutation, workspace, indexing, and termination behavior explicit.
\item Test a simple exact case, a difficult valid case, a boundary case, and an expected failure.
\end{itemize}}
\end{cornell}
\begin{summarybox}{Section summary}
Dawson's integral admits stable local series and asymptotic behavior, providing a scaled alternative to exponentially large intermediate expressions.  Select this technique only when its assumptions match the data, and accept a result only after an independent residual, error estimate, invariant, or refinement check.
\end{summarybox}

\section{6.10 Generalized Fermi-Dirac Integrals}
\begin{cornell}
\noterow{What is the central idea?}{Parameter-dependent integrals require reliable quadrature, transformations, and asymptotics across regimes where the integrand changes concentration.}
\noterow{How should it be used?}{For Generalized Fermi-Dirac Integrals, begin from explicit inputs, outputs, preconditions, and representation choices.  Classify the argument region and requested scaling; use identities to reach a stable region; choose series, recurrence, continued fraction, asymptotic expansion, or path continuation; cross-check identities.}
\noterow{Quantitative lens}{Identify the governing equation, recurrence, probability law, or geometric predicate for Generalized Fermi-Dirac Integrals.  Define every variable and scale; then derive a residual, error estimate, or invariant that can be checked independently.}
\noterow{Cost and reliability}{Analyze arithmetic work, storage, data movement, and convergence separately.  Relevant failure pressures include overflow before cancellation, wrong recurrence direction, branch errors, loss of tail probability, discontinuous region switching, and inaccurate normalization.  Use scaled tolerances and report both success criteria and diagnostic evidence.}
\noterow{Implementation checklist}{\begin{itemize}
\item Validate dimensions, domains, ordering, and structural assumptions.
\item Guard small divisors, invalid states, overflow, underflow, and nonfinite values.
\item Make mutation, workspace, indexing, and termination behavior explicit.
\item Test a simple exact case, a difficult valid case, a boundary case, and an expected failure.
\end{itemize}}
\end{cornell}
\begin{summarybox}{Section summary}
Parameter-dependent integrals require reliable quadrature, transformations, and asymptotics across regimes where the integrand changes concentration.  Select this technique only when its assumptions match the data, and accept a result only after an independent residual, error estimate, invariant, or refinement check.
\end{summarybox}

\section{6.11 Inverse of the Function x log(x)}
\begin{cornell}
\noterow{What is the central idea?}{Inversion of a monotone transcendental relation uses bracketing or transformed Newton iterations with branch-aware initial estimates.}
\noterow{How should it be used?}{For Inverse of the Function x log(x), begin from explicit inputs, outputs, preconditions, and representation choices.  Classify the argument region and requested scaling; use identities to reach a stable region; choose series, recurrence, continued fraction, asymptotic expansion, or path continuation; cross-check identities.}
\noterow{Quantitative lens}{Identify the governing equation, recurrence, probability law, or geometric predicate for Inverse of the Function x log(x).  Define every variable and scale; then derive a residual, error estimate, or invariant that can be checked independently.}
\noterow{Cost and reliability}{Analyze arithmetic work, storage, data movement, and convergence separately.  Relevant failure pressures include overflow before cancellation, wrong recurrence direction, branch errors, loss of tail probability, discontinuous region switching, and inaccurate normalization.  Use scaled tolerances and report both success criteria and diagnostic evidence.}
\noterow{Implementation checklist}{\begin{itemize}
\item Validate dimensions, domains, ordering, and structural assumptions.
\item Guard small divisors, invalid states, overflow, underflow, and nonfinite values.
\item Make mutation, workspace, indexing, and termination behavior explicit.
\item Test a simple exact case, a difficult valid case, a boundary case, and an expected failure.
\end{itemize}}
\end{cornell}
\begin{summarybox}{Section summary}
Inversion of a monotone transcendental relation uses bracketing or transformed Newton iterations with branch-aware initial estimates.  Select this technique only when its assumptions match the data, and accept a result only after an independent residual, error estimate, invariant, or refinement check.
\end{summarybox}

\section{6.12 Elliptic Integrals and Jacobian Elliptic Functions}
\begin{cornell}
\noterow{What is the central idea?}{Symmetric integral forms and duplication algorithms provide a unified, stable route to elliptic quantities and their inverses.}
\noterow{How should it be used?}{For Elliptic Integrals and Jacobian Elliptic Functions, begin from explicit inputs, outputs, preconditions, and representation choices.  Classify the argument region and requested scaling; use identities to reach a stable region; choose series, recurrence, continued fraction, asymptotic expansion, or path continuation; cross-check identities.}
\noterow{Quantitative lens}{Identify the governing equation, recurrence, probability law, or geometric predicate for Elliptic Integrals and Jacobian Elliptic Functions.  Define every variable and scale; then derive a residual, error estimate, or invariant that can be checked independently.}
\noterow{Cost and reliability}{Analyze arithmetic work, storage, data movement, and convergence separately.  Relevant failure pressures include overflow before cancellation, wrong recurrence direction, branch errors, loss of tail probability, discontinuous region switching, and inaccurate normalization.  Use scaled tolerances and report both success criteria and diagnostic evidence.}
\noterow{Implementation checklist}{\begin{itemize}
\item Validate dimensions, domains, ordering, and structural assumptions.
\item Guard small divisors, invalid states, overflow, underflow, and nonfinite values.
\item Make mutation, workspace, indexing, and termination behavior explicit.
\item Test a simple exact case, a difficult valid case, a boundary case, and an expected failure.
\end{itemize}}
\end{cornell}
\begin{summarybox}{Section summary}
Symmetric integral forms and duplication algorithms provide a unified, stable route to elliptic quantities and their inverses.  Select this technique only when its assumptions match the data, and accept a result only after an independent residual, error estimate, invariant, or refinement check.
\end{summarybox}

\section{6.13 Hypergeometric Functions}
\begin{cornell}
\noterow{What is the central idea?}{Analytic continuation and path integration extend a defining series while respecting singular points and branch choices.}
\noterow{How should it be used?}{For Hypergeometric Functions, begin from explicit inputs, outputs, preconditions, and representation choices.  Classify the argument region and requested scaling; use identities to reach a stable region; choose series, recurrence, continued fraction, asymptotic expansion, or path continuation; cross-check identities.}
\noterow{Quantitative lens}{Identify the governing equation, recurrence, probability law, or geometric predicate for Hypergeometric Functions.  Define every variable and scale; then derive a residual, error estimate, or invariant that can be checked independently.}
\noterow{Cost and reliability}{Analyze arithmetic work, storage, data movement, and convergence separately.  Relevant failure pressures include overflow before cancellation, wrong recurrence direction, branch errors, loss of tail probability, discontinuous region switching, and inaccurate normalization.  Use scaled tolerances and report both success criteria and diagnostic evidence.}
\noterow{Implementation checklist}{\begin{itemize}
\item Validate dimensions, domains, ordering, and structural assumptions.
\item Guard small divisors, invalid states, overflow, underflow, and nonfinite values.
\item Make mutation, workspace, indexing, and termination behavior explicit.
\item Test a simple exact case, a difficult valid case, a boundary case, and an expected failure.
\end{itemize}}
\end{cornell}
\begin{summarybox}{Section summary}
Analytic continuation and path integration extend a defining series while respecting singular points and branch choices.  Select this technique only when its assumptions match the data, and accept a result only after an independent residual, error estimate, invariant, or refinement check.
\end{summarybox}

\section{6.14 Statistical Functions}
\begin{cornell}
\noterow{What is the central idea?}{Distribution densities, cumulative probabilities, tails, and inverses should share consistent special-function kernels and tail-safe evaluations.}
\noterow{How should it be used?}{For Statistical Functions, begin from explicit inputs, outputs, preconditions, and representation choices.  Classify the argument region and requested scaling; use identities to reach a stable region; choose series, recurrence, continued fraction, asymptotic expansion, or path continuation; cross-check identities.}
\noterow{Quantitative lens}{Identify the governing equation, recurrence, probability law, or geometric predicate for Statistical Functions.  Define every variable and scale; then derive a residual, error estimate, or invariant that can be checked independently.}
\noterow{Cost and reliability}{Analyze arithmetic work, storage, data movement, and convergence separately.  Relevant failure pressures include overflow before cancellation, wrong recurrence direction, branch errors, loss of tail probability, discontinuous region switching, and inaccurate normalization.  Use scaled tolerances and report both success criteria and diagnostic evidence.}
\noterow{Implementation checklist}{\begin{itemize}
\item Validate dimensions, domains, ordering, and structural assumptions.
\item Guard small divisors, invalid states, overflow, underflow, and nonfinite values.
\item Make mutation, workspace, indexing, and termination behavior explicit.
\item Test a simple exact case, a difficult valid case, a boundary case, and an expected failure.
\end{itemize}}
\end{cornell}
\begin{summarybox}{Section summary}
Distribution densities, cumulative probabilities, tails, and inverses should share consistent special-function kernels and tail-safe evaluations.  Select this technique only when its assumptions match the data, and accept a result only after an independent residual, error estimate, invariant, or refinement check.
\end{summarybox}

\section*{Method comparison}
\begin{center}
\small
\begin{tabularx}{\linewidth}{>{\bfseries\raggedright\arraybackslash}p{0.22\linewidth}XX}
\toprule
Method or lens & Best use & Main caution \\
\midrule
Series & Small or central arguments & Stop from a rigorous term estimate \\
Continued fraction & Complementary tails & Use scaled recurrences and tiny-denominator guards \\
Asymptotic expansion & Large arguments & Optimal truncation may precede term growth \\
\bottomrule
\end{tabularx}
\end{center}

\section*{Worked example}
\begin{exambox}{Independent worked example}
Compute $\Gamma(6)$ from the recurrence $\Gamma(z+1)=z\Gamma(z)$ and $\Gamma(1)=1$: $\Gamma(6)=5\cdot4\cdot3\cdot2\cdot1=120$.  The recurrence is exact here.  For large noninteger arguments, evaluate $\log\Gamma$ to avoid overflow and recover ratios through log differences.
\end{exambox}

\section*{Chapter synthesis}
\begin{summarybox}{Chapter synthesis}
The chapter's methods should be viewed as a decision system rather than a list of routines.  Begin with the mathematical problem and its conditioning, exploit structure, select an algorithm with compatible stability and cost, and verify the computed answer with evidence that is independent of the internal iteration.  The recurring implementation discipline is: classify the argument region and requested scaling; use identities to reach a stable region; choose series, recurrence, continued fraction, asymptotic expansion, or path continuation; cross-check identities.
\end{summarybox}

\subsection*{Common mistakes and numerical hazards}
\begin{warnbox}{Common mistakes and numerical hazards}
Overflow before cancellation, wrong recurrence direction, branch errors, loss of tail probability, discontinuous region switching, and inaccurate normalization.  A second recurring mistake is reporting only a computed value: always retain tolerances, iteration counts, residuals, refinement behavior, and relevant structural checks.
\end{warnbox}

\subsection*{Retrieval practice}
\textbf{Questions}
\begin{enumerate}
\item What problem does Gamma Function, Beta Function, Factorials, Binomial Coefficients address, and which assumption is easiest to violate?
\item What problem does Incomplete Gamma Function and Error Function address, and which assumption is easiest to violate?
\item What problem does Exponential Integrals address, and which assumption is easiest to violate?
\item What problem does Hypergeometric Functions address, and which assumption is easiest to violate?
\item What problem does Statistical Functions address, and which assumption is easiest to violate?
\item How do conditioning, stability, and the final residual play different roles in this chapter?
\end{enumerate}

\textbf{Brief answers}
\begin{enumerate}
\item Logarithmic forms and identities extend factorial-like quantities while avoiding immediate overflow and preserving sign information. Recheck the section's domain, structure, scaling, and failure diagnostics.
\item Series and continued fractions cover complementary regions; selecting by argument prevents slow convergence and loss of accuracy. Recheck the section's domain, structure, scaling, and failure diagnostics.
\item Exponential integrals need different representations near singular points, across branch-sensitive regions, and for large arguments. Recheck the section's domain, structure, scaling, and failure diagnostics.
\item Analytic continuation and path integration extend a defining series while respecting singular points and branch choices. Recheck the section's domain, structure, scaling, and failure diagnostics.
\item Distribution densities, cumulative probabilities, tails, and inverses should share consistent special-function kernels and tail-safe evaluations. Recheck the section's domain, structure, scaling, and failure diagnostics.
\item Conditioning describes sensitivity of the mathematical problem, stability describes error propagation by the algorithm, and the residual measures satisfaction of the computed equations; none is a universal substitute for the others.
\end{enumerate}

\subsection*{Mastery checklist}
\begin{itemize}
\item[$\square$] I can explain and select Gamma Function, Beta Function, Factorials, Binomial Coefficients without relying on a memorized code listing.
\item[$\square$] I can explain and select Incomplete Gamma Function and Error Function without relying on a memorized code listing.
\item[$\square$] I can explain and select Exponential Integrals without relying on a memorized code listing.
\item[$\square$] I can explain and select Incomplete Beta Function without relying on a memorized code listing.
\item[$\square$] I can explain and select Bessel Functions of Integer Order without relying on a memorized code listing.
\item[$\square$] I can state a scaled stopping rule and an independent verification check.
\item[$\square$] I can identify at least one conditioning risk and one algorithmic stability risk.
\item[$\square$] I can estimate time, storage, and data-movement costs at the appropriate level.
\end{itemize}

\end{document}

